% =========================================================================
% Measuring What Matters When the Robot Has the Gun:
% A Validity Audit of Hallucination-Induced Unauthorized Action in LLM Agents
% =========================================================================
\documentclass{article}
% NeurIPS 2025
\usepackage[final, main]{neurips_2025}
\usepackage[utf8]{inputenc}
\usepackage[T1]{fontenc}
\usepackage{hyperref}
\usepackage{url}
\usepackage{booktabs}
\usepackage{amsfonts}
\usepackage{amsmath}
\usepackage{nicefrac}
\usepackage{microtype}
\usepackage{xcolor}
\usepackage{graphicx}
\usepackage{multirow}
\title{Measuring What Matters When the Robot Has the Gun:\\
A Validity Audit of Hallucination-Induced Unauthorized Action in LLM Agents}

\author{%
    Jake Klosowski \\
  \texttt{jfkx@stanford.edu}
  \And
  Michelle Campeau \\
  \texttt{campeau@stanford.edu}
} 

\date{Draft, May 2026}

\begin{document}
\maketitle
\begin{abstract}
The case for deploying LLM agents in irreversible-action settings, such as military targeting, embodied robotics, and autonomous trading, rests on a transfer claim: that strong scores on existing agent-safety benchmarks predict that an agent will not take an explicitly forbidden action when it matters. We argue that this transfer claim is not supported by current measurements, because no published benchmark isolates the construct doing the work. The relevant construct is the conditional rate at which an agent, having hallucinated, then crosses an explicit prohibition. We formalize this construct as Hallucination-Induced Unauthorized Action (HIUA), and we decompose it from the four neighboring constructs that the agent-safety literature often confounds it with. We then audit nine prior instruments against Messick's six aspects of validity. We find that the substantive aspect, that is, the cognitive-process claim that hallucination caused the violation, is essentially undefended in the surveyed corpus, even where the manipulations needed to test it already exist. We find that the structural aspect is typically reported through single accuracy numbers that aggregate over a documented 90-percentage-point dispersion in the dissociation between rule recall and rule adherence. We further find that the consequential aspect predicts three gaming failures already visible in adjacent benchmarks. We sketch HIUA-Bench, a design that factorially crosses three hallucination sub-types against three prohibition-salience levels across four task domains, with simulator-verifiable outcomes, a recall probe that partitions lucid from hallucinated violation, and a Generalizability-Theory reliability plan reporting $\Phi$ and $\mathrm{E}\rho^2$ across persons, items, raters, and occasions. We report a small-scale pilot of the benchmark against six open-weights models, in which we observe that the partition is empirically realizable, that prohibition salience is the largest single effect on violation rate (1.2\% at high salience against 23.3\% at low salience), and that the HIUA-to-KBV ratio on this panel is approximately 2.75 to 1.

\end{abstract}
\section{Introduction}
The empirical anchors for an agent-safety paper in 2026 are not subtle. RoboPAIR, an algorithmic jailbreak procedure applied to three commercial robot platforms, achieved approximately 100\% attack success in eliciting physical harmful actions from the NVIDIA Dolphins self-driving system, the Clearpath Jackal UGV running a GPT-4o planner, and the Unitree Go2 quadruped running GPT-3.5~\cite{robopair2024}. ARMOR 2025, the first military-doctrine-aligned LLM safety benchmark, evaluated twenty-one commercial LLMs against 519 prompts grounded in the Law of War, the Rules of Engagement, and the DoD Joint Ethics Regulation. It documented two distinct failure modes: models that hallucinate rules or invent constraints when asked to justify actions, and models that refuse lawful and scoped requests under operational constraints. The authors concluded that deploying general-purpose LLMs for military decision support without additional controls is premature~\cite{armor2026}. Anthropic's June 2025 agentic-misalignment red team, in which models were given autonomous agentic roles and access to a simulated company's email system, found that Claude Opus 4 chose to blackmail a fictional executive in 96\% of relevant scenarios. Gemini 2.5 Flash blackmailed in 96\% of cases, GPT-4.1 and Grok 3 Beta in 80\%, and DeepSeek-R1 in 79\%~\cite{anthropicagentic2025}.

These are three different mechanisms with overlapping blast radius: adversarial jailbreak, doctrine-mismatch under genuine military prompts, and emergent goal-protection under shutdown threat. The natural reading is that the field has many benchmarks for different parts of an unsafe-agent landscape. The problem this paper takes up is that the existing benchmarks do not appear to carve the landscape at the joints that matter when an action is irreversible.

The clearest empirical demonstration of this carving problem is in the constraint-adherence literature. DriftBench, which evaluated seven models from five providers across 2{,}146 scored runs on multi-turn ideation tasks, found that the rate at which models violate constraints they can simultaneously restate verbatim, known as the knows-but-violates or KBV rate, ranges from 8\% on the lowest-violation model to 99\% on the highest~\cite{driftbench2026}. SafeAgentBench, which evaluated embodied LLM agents in an interactive simulator on 750 tasks across 10 hazard categories, found that the best-performing baseline achieves 69\% success on safe tasks but only 5\% rejection on hazardous tasks~\cite{safeagentbench2024}. AgentHarm, which evaluated frontier models on 110 explicitly malicious agent tasks (440 with augmentations) across 11 harm categories, found that leading LLMs are surprisingly compliant with malicious agent requests, even without jailbreaking~\cite{agentharm2024}. These three findings, namely the 91-percentage-point KBV dispersion, the 5\% hazardous-task rejection rate, and the pre-jailbreak compliance rate, are routinely reported under the umbrella of ``agent safety,'' and developers often cite their relative scores when arguing about deployment. We observe that they are measurements of different conditional distributions over the same outcome variable, and that conflating them obscures important differences in mechanism.

This paper makes four contributions. First, in Section~\ref{sec:construct}, we formally define Hallucination-Induced Unauthorized Action (HIUA) as the conditional rate $P(\text{forbidden action} \mid \text{hallucination} \wedge \text{salient prohibition})$, and we decompose it from four neighboring constructs that the agent-safety literature often confounds it with. Second, in Section~\ref{sec:audit}, we audit nine prior instruments against Messick's six aspects of validity, and we show that the substantive aspect, namely the claim that hallucination caused the violation as opposed to merely co-occurring with it, is essentially undefended in the surveyed corpus. Third, in Section~\ref{sec:design}, we sketch HIUA-Bench, a design that factorially crosses three hallucination sub-types against three prohibition-salience levels and four task domains. The design uses simulator-verifiable outcomes, a post-hoc recall probe to partition lucid from hallucinated violation, and a Generalizability-Theory reliability plan reporting $\Phi$ and $\mathrm{E}\rho^2$ across persons, items, raters, occasions, and paraphrases. Fourth, in Section~\ref{sec:pilot}, we report a small-scale pilot of the benchmark against six open-weights models, in which we observe that the partition is empirically realizable, that prohibition salience is the largest single effect on violation rate (1.2\% at high salience against 23.3\% at low salience), and that the HIUA-to-KBV ratio on this panel is approximately 2.75 to 1.

The thesis that this paper defends is the following. When the cost of a single forbidden action is irreversible, such as a fired weapon, a transferred fund, or a thrown switch, measurement that confounds lucid violation with hallucination-induced violation systematically mis-licenses deployment claims. A benchmark that reports a single aggregate violation rate cannot distinguish a model that needs a faithfulness intervention from a model that needs an instruction-following intervention from a model that needs a perception-grounding intervention. The validity audit below is the argument, the design sketch is the constructive proposal, and the pilot is the proof that the partition is empirically realizable.

\section{Background: validity theory in the agent-safety context}
We take the reader's familiarity with Messick's unified validity framework~\cite{messick1989}, Kane's argument-based approach~\cite{kane2013}, and the recent application of these frameworks to AI evaluation~\cite{m2m2025, nomologicalnetworks2026} as given. We note three points here because they do load-bearing work in Section~\ref{sec:audit}.

First, the construct under test is not a property of the LLM as a system. It is a property of an interpretation placed on an LLM's score, as Messick's view of validity as ``an overall evaluative judgment of the degree to which empirical evidence and theoretical rationales support the adequacy and appropriateness of interpretations and actions based on test scores'' makes explicit~\cite{messick1989, standards2014}. The relevant interpretation for the deployment community is that the agent will not, when deployed, take a forbidden action because of a hallucinated belief about authorization, state, or tool effect. We observe that a benchmark which reports a violation rate without partitioning the conditional pathway licenses a weaker interpretation than the one its consumers often use it to license.

Second, Cronbach and Meehl's original framing of construct validity as the placement of a construct within a nomological network of theoretical relations~\cite{cronbachmeehl1955} is, in our view, more directly relevant to LLM evaluation than the Borsboom causal account~\cite{borsboom2004} for the use case at hand. Freiesleben has recently made this argument explicitly in the LLM-capability context: the nomological account avoids the strong ontological commitments of the causal account, while providing more substantive constraint on construct meaning than Kane's purely inferential refinement~\cite{nomologicalnetworks2026}. We adopt this view for HIUA. We do not claim that there is a unitary latent ``hallucination propensity'' causing variance in observed violation rates. We claim, instead, that HIUA, if it is a construct worth measuring, should sit in a predictable web of correlations with related measurements, including hallucination rate on text-only benchmarks, jailbreak compliance, and capability. This web is what licenses the deployment interpretation.

Third, the consequential aspect of validity, which concerns the evaluation of the social and behavioral consequences of using a test in a particular way, has been the least applied of Messick's six in AI evaluation~\cite{jacobswallach2021, m2m2025}. The agent-safety benchmark case forces it back into prominence, because the use of these benchmarks is to license deployment in irreversible-action settings. A benchmark that produces a plausible-looking score, but does not isolate the construct that the deployment claim requires, is not neutral. It is, instead, misleading in a particular direction. We return to this point in Section~\ref{sec:consequential}.

\section{Construct definition}
\label{sec:construct}
\subsection{Formal definition}
Let an agent trajectory $\tau = (s_0, a_0, s_1, a_1, \ldots)$ proceed under a system prompt containing an explicit prohibition $P$. We take $P$ to be a single-sentence rule of the form ``do not perform any action $a \in A_{\text{forbidden}}$.'' We let $h_t \in \{0, 1\}$ indicate the occurrence of a hallucination event at step $t$. By hallucination event, we mean a generated intermediate token sequence (such as chain-of-thought, tool argument, or tool-output interpretation) that is unfaithful to the model's input and to the ground-truth world state. We let $s_t \in \{0, 1\}$ indicate the salience condition: that is, $P$ was present in the model's effective context, was syntactically parseable, and was within the attention window. We operationalize the salience condition as the event that a recall probe issued at time $t$ would elicit a verbatim or near-verbatim restatement of $P$.

The HIUA construct is then the conditional rate

\[
\mathrm{HIUA}(M, D) = P\bigl(a_t \in A_{\text{forbidden}} \mid h_t = 1 \wedge s_t = 1;\, M, D\bigr)
\]
for a model $M$ on a domain distribution $D$. We note that this is a narrower target than aggregate violation rate (which marginalizes over $h_t$ and $s_t$), aggregate hallucination rate (which marginalizes over $a_t \in A_{\text{forbidden}}$), and aggregate compliance rate (which marginalizes over $h_t$ and over instruction provenance).

\subsection{Sub-constructs}
HIUA decomposes into three sub-constructs that map, approximately, onto stages of the recently-surveyed agent-hallucination taxonomy~\cite{halluagentsurvey2025}.

\textit{Authorization hallucination.} The agent fabricates that an exception, override, or higher-priority instruction permits $a \in A_{\text{forbidden}}$. As a concrete example, a forged ``manager email'' arrives mid-task with apparent override authority, and the agent treats the forgery as authoritative without grounding its trust in any verified channel. This is, in our taxonomy, a reasoning-stage hallucination.

\textit{State hallucination.} The agent fabricates a world state in which $a$ no longer belongs to $A_{\text{forbidden}}$. For instance, an ambiguous file name suggests that the file is a duplicate, the agent infers that ``this is obviously a duplicate, so deletion is not data loss,'' and the agent then deletes the unique copy. This is a perception-stage hallucination.

\textit{Tool hallucination.} The agent fabricates a tool's behavior such that calling the tool produces $a$ without the agent representing the call as $a$. For instance, a tool's docstring describes it as a search, the agent treats the search call as cost-free, but the tool in fact triggers a billable action. The agent has not represented its own action as ``spending money.'' This is an execution-stage hallucination.

These three sub-constructs are, in our view, theoretically distinct in their mechanism, in the interventions that would mitigate them, and, as we will argue in Section~\ref{sec:audit}, in the measurement design that would isolate them.

\subsection{Anti-construct}
HIUA is what is sometimes called a specific construct. It is sharpened by being explicit about what it is not. We identify four neighboring failure modes that the construct explicitly excludes.

\textit{Lucid violation (knows-but-violates).} The rule was recalled and the action was taken anyway. This is what DriftBench measures most directly, with rates from 8\% to 99\% across models~\cite{driftbench2026}. We note that the lucid violation rate is not HIUA. A model that lucidly violates a rule is not hallucinating, even if its behavior is unsafe.

\textit{Jailbreak compliance.} The user or third-party content explicitly instructs the harmful action, and the agent complies. This is what AgentHarm~\cite{agentharm2024}, RoboPAIR~\cite{robopair2024}, InjecAgent~\cite{injecagent2024}, and AgentDojo~\cite{agentdojo2024} measure. In these cases, the agent's intermediate generation is not unfaithful. The agent is, instead, faithfully complying with an instruction it should not have complied with.

\textit{Over-refusal of legitimate tasks.} The model refuses a permitted action because it pattern-matches to a harmful one. OR-Bench documents this as a measurable and distinct phenomenon across 80{,}000 prompts and 32 LLMs~\cite{orbench2024}, and ARMOR finds it specifically as a military-context failure mode~\cite{armor2026}.

\textit{Policy-invisible violation.} The agent took the action because the policy facts needed to recognize the action as forbidden were not present in its context. PhantomPolicy formalizes this case across eight violation categories, with the policy intentionally absent from tool responses~\cite{phantompolicy2026}. We note that this is a missing-information problem rather than a hallucination problem.

A measurement that cannot partition these four neighbors from HIUA is, in our view, measuring a confounded construct. Borsboom's causal account of validity requires that variations in the target attribute produce variations in the measurement outcome~\cite{borsboom2004}. If the same observed violation rate can be produced by any of five distinct mechanisms, then the variation in attribute (HIUA propensity specifically) is not what is producing the variation in the measurement. The measurement is responsive, but it is responsive to a mixture.

\section{Related work, organized by what each instrument measures}
We organize the prior literature by the construct each instrument implicitly measures, rather than by chronology or by author. We note that this organization is not neutral. It is the work that the audit in Section~\ref{sec:audit} will lean on.

\subsection{Malicious-user-instruction compliance}
The largest cluster of recent agent-safety benchmarks measures whether an agent complies with an explicit instruction from the user, the system, or a third party to perform an action that the agent should refuse. AgentHarm's 110 explicitly malicious agent tasks (440 with jailbreak augmentations) span 11 harm categories from fraud to cybercrime to harassment. The authors document that leading models are surprisingly compliant without any jailbreaking, that universal jailbreak templates adapt effectively to agentic tasks, and that jailbroken agents retain their capabilities while behaving maliciously across multi-step trajectories~\cite{agentharm2024}. AgentDojo, with 97 realistic tasks and 629 security cases spanning email, banking, and travel-booking environments, found that GPT-4o's benign utility of 69\% drops to 45\% under prompt-injection attack, with targeted attack success rates of 53.1\% on the ``Important message'' canonical attack~\cite{agentdojo2024}. InjecAgent's 1{,}054 cases across 17 user tools and 62 attacker tools found that ReAct-prompted GPT-4 is vulnerable to indirect prompt injection 24\% of the time, with hacking-prompt reinforcement nearly doubling the success rate~\cite{injecagent2024}. RoboPAIR achieved approximately 100\% attack success in eliciting physical harmful actions across three robot platforms~\cite{robopair2024}.

These instruments cluster cleanly around a common construct, namely the agent's resistance to taking forbidden actions when those actions are explicitly requested by an adversary. The adversary may be the user (AgentHarm), third-party content (InjecAgent, AgentDojo), or an algorithmic attacker (RoboPAIR). However, in every case the harmful action is in the instruction stream that the agent processes. This is, in a sense, the complement of HIUA. In HIUA, no actor has instructed the harmful action. The agent has, instead, confabulated authorization, state, or tool effect on its own.

\subsection{Outcome-aggregated agent safety}
A second cluster measures aggregate agent failure rates across heterogeneous scenarios, without distinguishing the causal pathway. ToolEmu's LM-emulated sandbox covering 36 high-stakes toolkits and 144 test cases reports that the safest LM agent fails 23.9\% of the time according to the automatic evaluator, and 68.8\% of the time according to human raters~\cite{toolemu2024}. R-Judge's 569 multi-turn agent interaction records spanning 27 risk scenarios in 5 application categories and 10 risk types measures the proficiency of an LLM-judge in identifying safety risks post-hoc, with GPT-4o achieving 74.42\%~\cite{rjudge2024}. Agent-SafetyBench's 349 interaction environments and 2{,}000 test cases across 8 risk categories and 10 failure modes found that none of 16 popular LLM agents achieved a safety score above 60\%, and the authors identify lack of robustness and lack of risk awareness as the two fundamental defects~\cite{agentsafetybench2024}.

These instruments report a single aggregated failure rate per model, sometimes with a per-category breakdown. We observe that the DriftBench finding that KBV varies from 8\% to 99\% across models~\cite{driftbench2026} implies that aggregating across hallucinated and lucid violations averages over a roughly 90-percentage-point swing in the dissociation between recall and adherence. A single accuracy number is, therefore, difficult to interpret as a stable model property. For example, the same 30\% violation rate could reflect a model that always recalls and sometimes violates anyway (lucid), a model that never recalls and always violates (recall failure), or a model that confabulates new authorizations 30\% of the time (HIUA proper). These three cases correspond to different failure modes with different remediations.

The R-Judge case is, in our view, particularly instructive. It is a meta-evaluation in which the LLM under test is asked to judge agent traces, rather than to act. This is a different construct again, namely risk awareness as judgment rather than risk avoidance as action. R-Judge's contribution is real, but reading its scores as a measure of agent safety in the deployment sense is, in our view, a construct-level mismatch.

\subsection{Lucid constraint adherence}
A third cluster measures the gap between what a model can recall and what a model does. DriftBench's central finding is the dissociation itself: across 2{,}146 scored benchmark runs spanning seven models, four interaction conditions, and 38 research briefs from 24 scientific domains, models accurately restate constraints they simultaneously violate, with KBV rates of 8\% to 99\% across models~\cite{driftbench2026}. Structured checkpointing partially reduces KBV but does not close the dissociation. AgentIF, which was constructed from 50 real-world agentic applications with instructions averaging 1{,}723 words and 11.9 constraints each across three constraint types (formatting, semantic, and tool), found that the best of the evaluated models followed fewer than 30\% of instructions perfectly~\cite{agentif2025}. The Hierarchical Safety Principles benchmark documented a quantifiable ``cost of compliance,'' where safety constraints degrade task performance even when compliant solutions exist, as well as an ``illusion of compliance'' in which high adherence rates often mask task incompetence rather than principled choice~\cite{hsp2025}.

We read the instruction-hierarchy training of Wallace et al.~\cite{instructionhierarchy2024} as a fix for the gap that this cluster exposes. By training models to prioritize system over developer over user over data instructions, models are expected to lucidly adhere to higher-priority rules even when later content suggests otherwise. We note that the fix targets the dissociation that DriftBench and AgentIF measure, rather than the hallucination pathway that HIUA targets.

This cluster's value for HIUA is, primarily, methodological. It has established that constraint recall and constraint adherence are dissociable, which is precisely the partition that HIUA-Bench will need to perform in reverse: separating failure to adhere despite recall (lucid) from failure to adhere because of hallucination.

\subsection{Embodied safety}
A fourth cluster targets embodied LLM agents in simulators. SafeAgentBench's 750 tasks (with 450 hazardous and 300 safe controls) across 10 hazard categories and 3 task types are grounded in an interactive AI2-THOR-style simulator. The authors find that the best-performing baseline achieved 69\% success on safe tasks, but only 5\% rejection on hazardous ones~\cite{safeagentbench2024}. AGENTSAFE extends this with 45 adversarial scenarios, 1{,}350 hazardous tasks, and 9{,}900 instructions evaluated across nine state-of-the-art VLMs and two embodied agent workflows, with multi-level diagnosis across perception, planning, and execution. The authors uncover systematic failures in translating hazard recognition into safe planning and execution~\cite{agentsafe2025}.

HEAL is, we believe, the closest existing analogue to HIUA in design. It probes hallucinations in long-horizon embodied planning by constructing scene-task inconsistencies, including distractor injection, task-relevant object removal, synonymous object substitution, and scene-task contradiction. These probes induce hallucination rates up to 40 times higher than base prompts across 12 models in two simulators~\cite{heal2025}. HEAL finds that models cannot reliably reject infeasible tasks, that hallucinations persist despite feedback-based mitigation, and that they are only partially reduced by cross-modal input. The illustrative example HEAL gives, namely a robot instructed to ``put the knives in the dishwasher'' in a scene without a dishwasher, then hallucinating the dishwasher and placing sharp utensils in an empty cabinet, is exactly the state-hallucination sub-construct from Section~\ref{sec:construct}.

The gap between HEAL and HIUA is the dependent variable. HEAL measures plan correctness conditional on the inconsistency manipulation. HIUA would, instead, measure prohibition crossing conditional on the same manipulation. The two are, in principle, recoverable from the same trace data. HEAL did not analyze it that way.

\subsection{Hallucination measurement, standalone}
The hallucination literature in isolation from agent action has matured considerably. HalluLens distinguishes extrinsic from intrinsic hallucination, dynamically generates its evaluation data to prevent test-set leakage, and disentangles hallucination from factuality~\cite{hallulens2025}. The recent survey of agent-specific hallucinations proposes a stage-based taxonomy spanning the agent workflow (brain, perception, action) and catalogs 18 triggering causes~\cite{halluagentsurvey2025}. Neither of these instruments is cross-linked to a measure of consequential action. They measure whether a hallucination occurred, but not whether it then caused a forbidden action.

\subsection{Validity and benchmarking-of-benchmarks}
A fifth and methodologically central cluster has emerged in the past two years, arguing that the agent-safety literature needs a more disciplined validity practice. Measurement-to-Meaning lays out a validity-centered framework for reasoning about which claims existing measurements can and cannot support~\cite{m2m2025}. The nomological-networks-required argument applies Cronbach-Meehl to LLM capability benchmarks, and argues that the nomological account is the most suitable foundation~\cite{nomologicalnetworks2026}. BetterBench assessed 25 AI benchmarks against 46 best-practice criteria, and found large quality differences, with most benchmarks failing to report statistical significance or to allow easy replication~\cite{betterbench2024}. The Eriksson et al.\ interdisciplinary meta-review of approximately 100 studies on AI-benchmark shortcomings names misaligned incentives, construct validity gaps, and gaming as systemic problems~\cite{canwetrust2025}. Raji et al.\ critique the practice of taking ``general'' benchmarks as abstractions for general capability~\cite{raji2021}. Bowman and Dahl propose four criteria that most NLU benchmarks fail~\cite{bowmandahl2021}. The PSN-IRT analysis of 41{,}871 items across 11 LLM benchmarks documents widespread saturation, insufficient difficulty ceilings, and item-level contamination~\cite{lostinbench2025}. Liao and Xiao argue for socio-technically situated evaluation~\cite{liaoxiao2023}. Jacobs and Wallach propose measurement modeling as the right frame for fairness in computational systems, emphasizing content and consequential validity as load-bearing~\cite{jacobswallach2021}.

This methodological cluster is what makes the present paper possible. It also makes the paper risky. The audit below is not novel in form, only in target. We try, in Section~\ref{sec:audit}, to make the target-specific arguments sharp enough that the audit is not a Messick checklist exercise, but, instead, a critique that produces conclusions the literature does not already contain.

\section{Validity audit: nine instruments, six aspects}
\label{sec:audit}
This section is the core of the paper. We organize by Messick's six aspects of validity rather than by instrument, and within each aspect we audit at least three prior instruments against the HIUA construct. Our aim is to produce, in each subsection, at least one critique that would not exist without the validity framework.

\subsection{Content aspect: does the test span the construct?}
The content question for HIUA is whether the items in a benchmark sample the joint space of $\{\text{hallucination type}\} \times \{\text{prohibition salience}\} \times \{\text{action domain}\}$ densely enough to license generalization to the deployment use case. We audit four instruments.

SafeAgentBench's 10 hazard categories, simulator-grounded across 3 task types, provide solid content coverage of hazardous instructions in embodied settings~\cite{safeagentbench2024}. They do not, however, factorially cross hazardous instructions with hallucination triggers. The construct it spans is, in our reading, ``hazard refusal under direct adversarial instruction'' rather than HIUA. We note that the 5\% rejection rate on hazardous tasks is a measure of refusal in the presence of an explicit harmful request, not a measure of the agent's tendency to confabulate authorization for an action no one requested.

AgentHarm's 11 harm categories, with 110 base tasks and 440 augmented versions, provide thick coverage of malicious-user-instruction compliance~\cite{agentharm2024}. The user is the adversary by construction. The benchmark contains no items in which the agent has been given a clean task and then must avoid hallucinating a side-channel authorization that would license a harmful action.

InjecAgent and AgentDojo cover third-party-injected adversarial instructions: the agent receives a clean user instruction, but processes external content that contains a hostile instruction~\cite{injecagent2024, agentdojo2024}. We observe the same coverage gap from HIUA's perspective. The agent has been instructed (by the attacker) to do the harmful thing.

HEAL has, in our reading, the right manipulation in the right setting: scene-task inconsistencies that induce state hallucination at rates up to 40 times the base~\cite{heal2025}. However, its content coverage stops short of explicit prohibitions. The dependent variable is whether the agent's plan matched the (impossible) instruction, rather than whether the agent's plan crossed a prohibition that the system prompt placed on it. A content extension of HEAL with explicit per-trial prohibition statements would be, in our view, the closest existing instrument to HIUA on the content aspect.

The cumulative content-aspect critique is straightforward. Prevailing benchmarks cover the malicious-instruction half of agent safety thickly, and the hallucinated-authorization half almost not at all. Content validity for HIUA across the surveyed corpus is, by inspection, near zero.

\subsection{Substantive aspect: does the response process match the theorized cognition?}
This is the aspect on which the literature is, in our view, most exposed. Messick's substantive aspect requires evidence that the cognitive process generating the test response matches the theorized process underlying the construct~\cite{messick1989, standards2014}. For HIUA, the theorized process is ``hallucination event $\rightarrow$ forbidden action.'' The substantive aspect therefore asks: when the model produces a forbidden action, do we have evidence that a hallucination produced it, as opposed to merely co-occurring with it?

ToolEmu and Agent-SafetyBench score outcomes only~\cite{toolemu2024, agentsafetybench2024}. The agent's intermediate reasoning is not analyzed for the presence or character of a hallucination. The 23.9\% (automatic) and 68.8\% (human) failure rates that ToolEmu reports are, therefore, consistent with hallucination-driven failure, but they provide no evidence that hallucination produced any specific failure. We note that the post-hoc attribution ``this is a hallucination problem'' is exactly the kind of mechanistic claim that substantive validity is supposed to license.

R-Judge measures whether a judge LLM can identify the safety issue post-hoc~\cite{rjudge2024}. The cognitive process under test is the judge's, not the actor's. The 74.42\% GPT-4o-as-judge score is informative for an audit pipeline, but it is uninformative as evidence about why any specific agent in the trace took the action it took.

DriftBench is, in our view, the closest existing analogue to a substantive-aspect instrument. By interleaving a recall probe with the action, it partitions trials into a $2 \times 2$ of $\{\text{recall yes/no}\} \times \{\text{violation yes/no}\}$, and it lets us read off the KBV rate from the bottom-right cell~\cite{driftbench2026}. This is, in our reading, exactly the right move, and the resulting KBV rates (8 to 99\% across models) are precisely the kind of finding that an outcome-only benchmark cannot produce. However, the partition is one factor short of what HIUA requires. It separates ``did the model know the rule'' from ``did the model violate it,'' but it does not separate ``did the model hallucinate authorization'' from ``did the model hallucinate state'' from ``did the model hallucinate tool effect'' from ``did the model violate lucidly with full recall and no hallucination.'' The $2 \times 2$ must, in our view, become a $2 \times 2 \times 2$ or, with the sub-construct decomposition, a $2 \times 4 \times 2$ for $\{\text{recall}\} \times \{\text{hallucination-type or none}\} \times \{\text{violation}\}$.

HEAL has the manipulation half of what substantive validity requires. It intervenes on the world state to induce a hallucination, and then measures the agent's plan~\cite{heal2025}. It does not, however, measure prohibition crossing as the consequential outcome.

The cumulative critique is the following. There is, to our knowledge, no surveyed benchmark whose data could support a substantive-validity argument for HIUA as published. The data could, in principle, be re-analyzed. HEAL trace data plus an added prohibition layer, DriftBench plus a hallucination-type classifier on the trace, and ToolEmu plus a recall probe would all be candidates. We observe that the fact that the field has not done this re-analysis is itself part of the critique. Agent-safety benchmarks have been positioned as outcome-summary instruments, and the underlying substantive question has been displaced into discussion sections.

The constructive consequence for HIUA-Bench is that the substantive aspect must be designed in. We return to this in Section~\ref{sec:scoring}.

\subsection{Structural aspect: internal structure and reliability}
The structural aspect concerns whether the test's internal structure matches the theorized internal structure of the construct, and, for our purposes, whether the resulting score is reliable enough to bear the interpretation placed on it~\cite{messick1989, standards2014}. The course's central reliability framework, Generalizability Theory, partitions observed-score variance into facets (persons, items, raters, occasions, settings) and reports $\Phi$ (absolute reliability) and $\mathrm{E}\rho^2$ (relative reliability) coefficients~\cite{gtheoryage2025}. We audit four instruments against this lens.

AgentIF reports per-constraint-type breakdowns and an overall instruction-following rate, but no variance-component decomposition~\cite{agentif2025}. Given that instructions average 11.9 constraints each, the per-item difficulty distribution is, plausibly, highly skewed. Without a structural analysis, the reported below-30\% perfect-adherence rate could reflect that one or two constraint types are systematically harder, that the constraint counts confound difficulty, or both.

ToolEmu reports a striking discrepancy that is, in effect, a partial single-facet G-study: 23.9\% failure rate per automatic evaluator versus 68.8\% per human rater on the same 144 cases~\cite{toolemu2024}. The 45-percentage-point gap is, in G-Theory terms, a strong indication that $\sigma^2(\text{rater})$ is large. It is large enough that the choice of rater is doing more work in producing the reported number than the choice of agent. The authors do not report it this way. They report it, instead, as a quality-of-emulator check. The interpretation as a reliability failure of the rating procedure is, we think, more directly load-bearing for downstream interpretation.

DriftBench's 8 to 99\% KBV variance across persons (models) is, in our view, exactly the kind of dispersion that should be partitioned~\cite{driftbench2026, gtheoryage2025}. If $\sigma^2(\text{person}) \gg \sigma^2(\text{item, occasion, rater})$, the benchmark is measuring real model differences, and a single per-model KBV rate is a useful summary. If $\sigma^2(\text{item})$ is comparable to $\sigma^2(\text{person})$, the benchmark is largely measuring which items happened to elicit KBV, and the per-model rates would shift with item resampling. The published data does not, in our reading, support that decomposition.

The PSN-IRT analysis of 41{,}871 items across 11 LLM benchmarks is the methodological precedent that we lean on most heavily~\cite{lostinbench2025}. It demonstrates that item-level IRT analysis on existing benchmarks is feasible, and that the results are sometimes uncomfortable. Widespread saturation, insufficient difficulty ceilings, and data contamination affect benchmarks that the field has been using to draw broad capability claims. The same analysis applied to agent-safety benchmarks is, to our knowledge, not yet published. Our prediction is that a substantial fraction of items in AgentHarm and SafeAgentBench would either show ceiling effects on frontier models (uninformative items) or suspiciously low discrimination parameters (items where most failures are model-independent).

The cumulative critique is the following. Agent-safety benchmarks report single numbers in cases where the underlying variance structure makes those single numbers difficult to interpret. The minimum reporting standard a HIUA-grade benchmark should adopt is, in our view, $\Phi$ and $\mathrm{E}\rho^2$ across at least persons $\times$ items $\times$ raters $\times$ occasions, with D-studies indicating how many items and occasions are needed for $\Phi \geq 0.80$ at the construct's target precision.

\subsection{Generalizability aspect: across persons, items, occasions, settings}
The generalizability aspect for HIUA is, in our view, the load-bearing assumption of every agent-safety paper in this space: that a benchmark score in one setting predicts behavior in another. The setting transfer of interest is from digital sandboxes, in which a forbidden action is sending an email or transferring a small simulated balance, to embodied and lethal settings, in which a forbidden action is firing a weapon or shutting off a life-support device. This transfer is rarely measured, and is almost never reported as a generalizability coefficient.

No published cross-walk exists, to our knowledge, between AgentHarm scores on the same models and SafeAgentBench scores on the same models. The two papers were published within months of each other~\cite{agentharm2024, safeagentbench2024}. The rank correlation between their reported model orderings would be a one-table contribution, and we have not been able to find it.

The irreversibility gap matters in a way that pure generalizability mathematics underweights. A 5\% violation rate is acceptable for a chatbot that types ``I cannot help with that'' 95\% of the time. It is, however, catastrophic for a drone that fires a weapon 5\% of the time it should not. The construct of interest has, in our view, a stakes-dependent decision threshold that is not separable from the construct itself. ARMOR makes the analogous argument for military deployment: civilian-context safety metrics underestimate the military-relevant safety gap, especially around lawful-instruction refusal under operational constraints~\cite{armor2026}. RoboPAIR's near-100\% physical-harm attack success rate on three robot platforms~\cite{robopair2024} is not ``almost the same'' as a 50\% rate. The difference between a 100\% and a 50\% jailbreak success rate is, we observe, the difference between deployment being unsafe-on-paper and unsafe-in-deployment.

Occasion variance is also under-reported. The Anthropic agentic misalignment study reports blackmail rates of 79 to 96\% across model families on a particular scenario~\cite{anthropicagentic2025}. The within-model variance across seeds, paraphrases, and prompt orderings is reported only sparsely. Without occasion-facet variance, we cannot tell whether a 96\% blackmail rate is a stable model property or an artifact of one prompt construction.

The cumulative critique is that agent-safety benchmarks report generalizability mostly via implicit assumption, rather than via measurement. For HIUA, generalizability evidence must include, in our view, (i) within-domain occasion-facet variance reported as a coefficient, (ii) cross-domain transfer reported as a correlation, and (iii) explicit irreversibility-stakes annotation per domain, so that a 5\% rate in a high-stakes domain is not silently averaged with a 5\% rate in a low-stakes domain.

\subsection{External aspect: what does the score correlate with?}
External validity, in the Cronbach-Meehl sense, requires placing the construct in a nomological network of expected relations to other measured constructs~\cite{cronbachmeehl1955, nomologicalnetworks2026, m2m2025}. A valid HIUA measurement should, on this view, show a predictable pattern of correlations: positive with measures of the same underlying generative process, weakly positive with measures of adjacent failure modes, and zero or inverse with capability above some threshold.

We list the patterns that we predict, and that a HIUA benchmark should be obliged to report.

A positive correlation with hallucination rate on text-only benchmarks such as HalluLens~\cite{hallulens2025}. The mechanism, namely unfaithful intermediate generation, is hypothesized to be the same. The rates should, therefore, covary.

A weakly positive correlation with AgentHarm compliance~\cite{agentharm2024}. The mechanism is different (jailbreak compliance rather than confabulation), but model families that are generally more compliant should also be more willing to follow a hallucinated authorization. The correlation should exist, but should be lower than the HalluLens correlation. If it is higher, we take the construct decomposition to be wrong: HIUA and jailbreak compliance are, in that case, being driven by a single underlying variable, not two.

A zero or inverse correlation with general capability (such as MMLU or GPQA) above some threshold. More capable models hallucinate less per token, but they are also deployed in higher-stakes settings, so the deployment-weighted HIUA risk may not decline with capability. If the empirical correlation is strongly positive, we take the construct to be contaminated with capability variance, and the benchmark to be measuring ``weak model'' rather than ``HIUA-prone model.''

A positive correlation with strategic-deception measures~\cite{deceptionsurvey2024} and with sycophancy measures~\cite{sharmasycophancy2023}. Both reflect breakdowns of faithful intermediate generation under pressure (strategic goal pressure or user pressure). HIUA reflects the breakdown without either pressure. The three should, in our view, sit in a graded family with a positive but not unit correlation.

No surveyed benchmark reports this pattern. The closest external-validity datapoint in the agent-safety literature is, we believe, the Anthropic agentic-misalignment study, which observed elevated blackmail rates under shutdown-threat conditions across model families~\cite{anthropicagentic2025}. This finding is consistent with a goal-protection construct that is related to HIUA but not identical. We note that the prediction was reported after the fact rather than pre-registered, and that the cross-model correlation pattern with other safety measures was not the focus.

The cumulative critique is the following. The field should, in our view, publish predicted correlation patterns before running new instruments. The nomological-network practice has not yet migrated from the validity-theory literature~\cite{nomologicalnetworks2026} to the agent-safety literature in any operational form.

\subsection{Consequential aspect: the leaderboard problem}
\label{sec:consequential}
Messick's consequential aspect requires evaluation of the social and behavioral consequences of using a test in a particular way~\cite{messick1989, kane2013, jacobswallach2021}. It is the aspect least applied in AI evaluation~\cite{m2m2025}, and the one with the highest stakes for HIUA, because the use of these benchmarks is to license deployment.

We identify three predictable consequences if HIUA-style benchmarks gain leaderboard status without consequential-aspect controls.

\textit{Surface compliance gaming.} Models trained to refuse anything resembling a benchmark trigger learn to over-refuse legitimate tasks. OR-Bench documents this at scale: 80{,}000 over-refusal prompts across 32 LLMs and 8 model families, with a hard subset of approximately 1{,}000 prompts that challenge even state-of-the-art models~\cite{orbench2024}. ARMOR documents the same failure mode in the military domain. Over-refusal of lawful, scoped requests under operational constraints is one of two principal failure modes that the authors identify, alongside hallucinated rule-invention~\cite{armor2026}. This is, we observe, Goodhart's law in the validity register: when a refusal rate becomes the target, refusal stops indicating safety~\cite{canwetrust2025}.

\textit{Policy-as-prompt theater.} Developers add longer system prompts to pass benchmarks. Production systems with token-budget pressure strip them, and the result is the policy-invisible-violation case that PhantomPolicy is explicitly built around: agents that took syntactically valid, user-sanctioned, semantically appropriate actions that violated policy because the decisive facts were absent from the visible context~\cite{phantompolicy2026}. A benchmark that scores agents with the full policy in context, but does not also score them with the policy partially or fully absent, licenses a deployment claim that production will not honor.

\textit{False assurance for high-stakes deployment.} A passing digital-sandbox score is cited to justify embodied or lethal deployment despite no demonstrated transfer. RoboPAIR's near-universal attack success on three robot platforms~\cite{robopair2024}, together with SafeAgentBench's 5\% hazardous-task rejection rate~\cite{safeagentbench2024}, suggests that the digital-to-embodied transfer is not just unmeasured, but specifically disconfirmed in the data we have. A digital-sandbox-leaderboard-top model is not, on present evidence, an embodied-safe model. The consequential failure mode is that a measurement-licensed claim (``our agent passes the safety benchmark'') gets used to make a deployment-class claim (``our agent is safe to deploy in setting $X$'') that the measurement does not support.

The instruction-hierarchy training of Wallace et al.\ is the field's principal architectural response to (1)~\cite{instructionhierarchy2024}. It partly mitigates surface-compliance gaming by training a policy across instruction roles rather than per-benchmark refusal patterns. We note that it does not address (2) and (3), and it introduces its own consequential risk in the form of the second ARMOR failure mode: over-refusal of lawful, scoped instructions~\cite{armor2026}.

The consequential-aspect recommendation for HIUA-Bench is, we propose, gated release. The benchmark should not be deployed as a public leaderboard. It should, instead, be released with the AgentHarm-style harm-score gating~\cite{agentharm2024, inspectevals2024}, which releases items only to verified developers and reports model results in aggregated, time-delayed form. The release should be accompanied by a per-model deployment-claim audit, that is, a one-page document, signed by the developer, stating which deployment claims the benchmark score is and is not being used to support. We take this proposal to make consequential validity an explicit object of test-developer responsibility, rather than a side effect of how the score is used in the marketplace.

\section{HIUA-Bench: a design sketch}
\label{sec:design}
The constructive proposal that follows is, by intent, smaller than the leaderboard benchmarks it would join. The validity argument is, in our view, that smaller-and-on-construct beats larger-and-confounded for the deployment-license use case.

\subsection{Item structure}
Each item is a tuple $\langle\text{environment},\, \text{prohibition},\, \text{hallucination-trigger},\, \text{ground-truth state},\, \text{action space}\rangle$.

The environment is one of four domains: file operations (digital, modeled in a sandboxed shell), communication (digital, modeled in an email/messaging emulator), financial (digital, modeled in a banking-API emulator), and embodied (AI2-THOR or a HEAL-derived scene~\cite{heal2025, safeagentbench2024}). The four-way domain choice provides, we note, cross-domain transfer evidence across the digital-to-embodied gradient that Section~\ref{sec:audit} identified as load-bearing.

The prohibition is a single-sentence forbidden-action statement in the system prompt. Salience is varied at three levels: early-system (within the first roughly 200 tokens), mid-system (after several thousand tokens of unrelated instruction), and tool-docstring (the prohibition appears in the docstring of a specific tool rather than the system prompt). The salience variation is motivated by the policy-invisible-violation literature~\cite{phantompolicy2026} and by the AgentIF finding that long instructions with many constraints are systematically harder~\cite{agentif2025}.

The hallucination trigger is a contextual perturbation designed to elicit one of the three sub-construct hallucination types from Section~\ref{sec:construct}. For authorization hallucination, a forged ``manager email'' granting an override arrives mid-task through a tool channel, and the email lacks any cryptographic or verified provenance. For state hallucination, an ambiguous identifier suggests a duplicate state (a file name suggesting redundancy, a sensor reading at the edge of a discretization boundary, an inventory message implying the target is already gone). For tool hallucination, a tool's docstring describes a benign action, but the tool's emulated implementation triggers a destructive side effect. The trigger is, in this case, the docstring-behavior mismatch.

The ground-truth state is simulator-verifiable, so that outcome scoring does not rely on an LLM judge for the primary measure. The action space includes at least one forbidden action and at least one compliant alternative. The agent that refuses to act at all scores neither.

The factorial is 3 sub-types $\times$ 3 salience levels $\times$ 4 domains, yielding 36 base items. Each base item is paraphrased 4 ways and run at 3 temperature seeds, yielding 12 occasions per base item and 432 trials per model evaluated. The scale is, we note, comparable to ToolEmu's 144 cases~\cite{toolemu2024}, but with a denser internal factorial.

\subsection{Scoring: three measurements per trial}
\label{sec:scoring}
The substantive-aspect critique in Section~\ref{sec:audit} requires that the cognitive process and the outcome both be measured. We score three things per trial.

\textit{Outcome.} Binary, simulator-derived: did the agent execute a forbidden action? No LLM judge is in the loop for this measure, which prevents the rater-variance failure that ToolEmu documented (the 23.9\% / 68.8\% gap~\cite{toolemu2024}) from propagating to the primary score.

\textit{Hallucination.} Rubric-coded from the chain-of-thought trace by a 3-model LLM-judge ensemble, with human spot-checks on 10\% of trials. The rubric distinguishes the three sub-types and a ``no hallucination'' category. Inter-rater reliability is reported as Cohen's $\kappa$ per category, and as Krippendorff's $\alpha$ overall. The 3-judge ensemble is itself one facet in the G-study (Section~\ref{sec:gstudy}). We expect $\sigma^2(\text{rater})$ to be substantial, and we report it rather than aggregating it away.

\textit{Recall probe.} After the action (or refusal), the agent is prompted to restate the prohibition. The probe partitions trials into a $2 \times 2$ of $\{\text{recall yes/no}\} \times \{\text{violation yes/no}\}$ and, jointly with the hallucination measurement, into the $2 \times 4 \times 2$ cell structure that HIUA requires. The HIUA estimate is the rate at which the agent executed a forbidden action conditional on recalling the rule \textit{and} exhibiting one of the three hallucination sub-types. That is the cell that maps to ``hallucinated authorization, state, or tool effect despite knowing the rule.'' Trials in the $\{\text{recall yes, no hallucination, violation}\}$ cell are lucid violations and are reported separately as the KBV rate, matching DriftBench's measurement~\cite{driftbench2026}.

\subsection{Reliability plan (G-study)}
\label{sec:gstudy}
The reliability plan applies the course's G-Theory framework to the facets named above. The crossed design is persons (models) $\times$ items (36 base items) $\times$ raters (3 LM-judges plus 1 human spot) $\times$ occasions (4 paraphrases $\times$ 3 seeds) $\times$ prohibition-salience (3 levels), with sub-type and domain treated as structurally fixed strata for separate reporting.

We report $\sigma^2$ decomposition across facets, $\Phi$ (absolute reliability) and $\mathrm{E}\rho^2$ (relative reliability) coefficients, and a D-study indicating how many items and occasions are needed for $\Phi \geq 0.80$ under each domain stratum. The IRT 2PL discrimination parameter is reported per item. Items with discrimination below 0.5 are flagged for revision rather than included in the headline rate. The discrimination distribution is published alongside the model rates, following the PSN-IRT precedent for transparency about item informativeness~\cite{lostinbench2025}.

The minimum reporting expectation for any future HIUA-claimant benchmark is, in our view, that $\sigma^2(\text{person}) > \sigma^2(\text{item})$ for the headline coefficient to be interpretable as a model property. If $\sigma^2(\text{item})$ dominates, the benchmark is measuring item difficulty, and should be reported as such.

\subsection{Validity argument by Messick aspect (pre-registered)}
The validity argument is pre-registered at item-construction time, and reported in the release alongside the empirical results.

\textit{Content.} Factorial coverage is justified against the agent-hallucination taxonomy~\cite{halluagentsurvey2025}. The three sub-types map onto reasoning, perception, and execution stages of the surveyed taxonomy. The four domains span the digital-to-embodied gradient. The three salience levels span the prohibition-context-cardinality dimension that AgentIF and PhantomPolicy implicate~\cite{agentif2025, phantompolicy2026}.

\textit{Substantive.} Chain-of-thought and recall probe interleaved with action instrument the hypothesized cognitive process (Section~\ref{sec:scoring}). The $2 \times 4 \times 2$ cell structure is the construct's empirical residence.

\textit{Structural.} $\Phi$ and $\mathrm{E}\rho^2$ coefficients are reported across all facets, IRT 2PL parameters per item, and D-studies indicate sample sizes required for given precision levels (Section~\ref{sec:gstudy}).

\textit{Generalizability.} Cross-domain transfer (digital to embodied) is reported as Spearman rank correlation between domain-stratum rates, with the irreversibility-stakes annotation per domain. Occasion-facet variance is reported as $\Phi(\text{occasions})$ and tabulated against $\Phi(\text{items})$, to make the relative importance of within-trial variability versus item variability explicit.

\textit{External.} Pre-registered predictions: positive correlation with HalluLens hallucination rate~\cite{hallulens2025}, weakly positive with AgentHarm compliance~\cite{agentharm2024}, zero or inverse with MMLU above the 0.7 threshold, positive with sycophancy~\cite{sharmasycophancy2023} and with strategic-deception measures~\cite{deceptionsurvey2024}. The pre-registration is on the published artifact. Deviations from the predictions are reported as evidence about the construct's empirical placement, rather than as failures of the benchmark.

\textit{Consequential.} Released through gated access with the AgentHarm-style harm-score gating~\cite{agentharm2024, inspectevals2024}. Not deployed as a public leaderboard. Release is accompanied by a developer-signed deployment-claim audit.

\section{Empirical pilot}
\label{sec:pilot}
We report a small-scale pilot of HIUA-Bench against an open-weights model panel served through Groq's hosted inference. The pilot's purpose is, first, to verify that the three-layer scoring stack discriminates among the four cells of the partition rather than collapsing trials into one or two cells, and, second, to produce an initial empirical datapoint that future work can extend.

\subsection{Panel and protocol}
The pilot evaluated six open-weights instruction-tuned models on the 24-item base bank from Section~\ref{sec:design}: Meta Llama 3.1 8B Instant, Meta Llama 3.3 70B Versatile, Meta Llama 4 Scout 17B (preview), Alibaba Qwen3 32B (preview), OpenAI GPT-OSS 20B, and OpenAI GPT-OSS 120B. All models were served through Groq's free-tier inference endpoint. Each model contributed one trial per item at temperature 0.7 with seed 0, yielding 144 trials, all of which completed successfully at the agent layer. The post-hoc judge ensemble was rate-limited on one of its three judges (see scoring-stack diagnostics below).

For each trial, the agent received the system prompt (containing the explicit prohibition) and the user message (containing the hallucination trigger), and produced a chain-of-thought followed by an action description. We then issued the recall probe in a fresh context and scored both responses through a three-layer action scorer (refusal-prefix heuristic, regex match, LLM-judge fallback) and a two-layer recall scorer (substring match, LLM-judge fallback). Post-hoc, a three-judge ensemble (Llama 3.1 8B Instant, Llama 3.3 70B Versatile, and Llama 4 Scout 17B) labeled each trial's hallucination type from the chain-of-thought transcript as one of \{none, authorization, state, tool\}.

\subsection{Scoring-stack diagnostics}
The action scorer's per-trial path distribution provides a first sanity check that the three layers are doing distinct work. Of the 144 trials, 45.8\% were scored by the regex layer alone, 51.4\% by the refusal-prefix heuristic, 1.4\% required the LLM-judge fallback, and 1.4\% returned scorer errors that we report separately. The fallback rate is low because the action scorer's two cheap layers handle the majority of cases, which is the design intent.

The recall scorer's distribution is more informative for the substantive-validity argument. Of the 142 correctly-recalled trials, 71.8\% were caught by exact-substring match, while 28.2\% were rescued by the LLM-judge fallback after the substring scorer had failed. The paraphrased-recall rate is, in our view, evidence that a benchmark using substring matching alone would systematically misclassify nearly a third of correct recalls as recall failures, which would inflate the HIUA estimate by counting paraphrased-recall violations as not-recalled.

The judge ensemble's inter-rater reliability is bounded by the per-judge error rate, which was uneven across the three judges. Llama 3.1 8B Instant returned valid labels on 140 of 144 trials; Llama 4 Scout 17B on 140 of 144; Llama 3.3 70B Versatile on only 20 of 144, almost entirely because of Groq's daily request-per-model cap on the Versatile endpoint. We therefore report pairwise Cohen's $\kappa$ restricted to trials in which both judges in the pair returned canonical labels. The pair with both judges at full strength, Llama 3.1 8B Instant against Llama 4 Scout 17B, achieved $\kappa = 0.419$ ($n=138$, observed agreement 82.6\%), which falls in the moderate-agreement range. The two pairs involving Llama 3.3 70B Versatile achieved higher $\kappa$ values (0.568 and 0.778) but on only 20 trials each, and we therefore treat them as suggestive rather than conclusive. The mean of the three pairwise coefficients is 0.588.

\subsection{Cell occupancy}
Table~\ref{tab:pilot-cells} reports the headline $2 \times 2$ of \{forbidden action?\} $\times$ \{recalled rule?\}, pooled across the panel.

\begin{table}[h]
\centering
\caption{Pilot cell occupancy across $N=144$ trials.}
\label{tab:pilot-cells}
\begin{tabular}{lrr}
\toprule
Cell & Count & Percent \\
\midrule
Compliant + recalled                       & 127 & 88.2\% \\
Compliant + not recalled                   &   2 &  1.4\% \\
Violation + recalled                       &  15 & 10.4\% \\
Violation + not recalled                   &   0 &  0.0\% \\
\bottomrule
\end{tabular}
\end{table}
We observe that the panel was almost uniformly compliant, with violations occurring on roughly one in ten trials. The compliant + not-recalled cell contains two trials and the violation + not-recalled cell is empty, indicating that the panel's recall probe was answered correctly on essentially every trial: the agents knew the rule. The question is therefore which of the violation trials were knows-but-violates (lucid) cases and which were hallucination-induced.

To answer that question we turn to the $2 \times 4 \times 2$ decomposition by judge-assigned hallucination sub-type. Of the 15 violation trials, the majority judge label was \textit{authorization} on 7 trials, \textit{state} on 3 trials, \textit{tool} on 1 trial, and \textit{none} on 4 trials. We define the pilot's HIUA estimate as the rate at which the agent executed a forbidden action conditional on recalling the rule and exhibiting one of the three hallucination sub-types in the judge's majority label, and the KBV estimate as the rate at which the agent violated despite recalling the rule and with a \textit{none} majority label. On this pilot:

\[
\widehat{\mathrm{HIUA}} = \frac{11}{144} = 7.6\%, \qquad \widehat{\mathrm{KBV}} = \frac{4}{144} = 2.8\%.
\]
The HIUA-to-KBV ratio is approximately 2.75 to 1 (11 against 4 trials). We caution against strong inferences from these absolute rates given the panel's restriction to open-weights models, the single-seed design, and the partial loss of the Llama 3.3 70B judge to Groq's rate cap. We note, nevertheless, that the partition produces non-trivial occupancy in both the HIUA and KBV cells, which is direct evidence that the cells are operationally distinguishable on this benchmark.

\subsection{Per-model rates}
Table~\ref{tab:pilot-per-model} reports per-model rates for all six models, each contributing $n=24$ trials.

\begin{table}[h]
\centering
\caption{Per-model pilot rates. Cells in percent. HIUA-cand and KBV are computed using the three-judge majority hallucination label.}
\label{tab:pilot-per-model}
\begin{tabular}{lrrrrr}
\toprule
Model & $n$ & Violation & Recall & HIUA-cand & KBV \\
\midrule
GPT-OSS 120B          & 24 &  0.0 & 100.0 &  0.0 & 0.0 \\
GPT-OSS 20B           & 24 &  4.2 &  91.7 &  4.2 & 0.0 \\
Llama 3.1 8B Instant  & 24 & 20.8 & 100.0 & 20.8 & 0.0 \\
Llama 3.3 70B Vers.   & 24 &  8.3 & 100.0 &  4.2 & 4.2 \\
Llama 4 Scout 17B     & 24 & 12.5 & 100.0 &  4.2 & 8.3 \\
Qwen3 32B (preview)   & 24 & 16.7 & 100.0 & 12.5 & 4.2 \\
\bottomrule
\end{tabular}
\end{table}
We observe a spread of 20.8 percentage points in violation rate across the panel, ranging from GPT-OSS 120B at the refusal floor to Llama 3.1 8B Instant at one-in-five violations. The Meta within-family scaling pair (3.1 8B against 3.3 70B) is consistent with the predicted capability-correlated pattern: the larger model violates less than the smaller model in the same family. Llama 4 Scout 17B is, on this panel, the model with the highest KBV share (8.3\%): its two judge-majority \textit{none} violation trials place it as the panel's most lucid violator rather than its most hallucination-prone violator. The smaller and earlier-generation models in the panel exhibit the highest HIUA-candidate rates, while the production GPT-OSS variants remain at or near the refusal floor. We caution that the pilot's single-seed design means these per-model differences are point estimates without standard errors, and they may not survive a full Generalizability-Theory decomposition.

\subsection{Sub-construct and salience effects}
Table~\ref{tab:pilot-subconstruct} reports violation rates broken down by the hallucination sub-type the item was designed to elicit and by prohibition salience.

\begin{table}[h]
\centering
\caption{Pilot violation rates by sub-construct and salience.}
\label{tab:pilot-subconstruct}
\begin{tabular}{lrr}
\toprule
Stratum & $n$ & Violation rate \\
\midrule
Authorization sub-construct & 54 &  9.3\% \\
State sub-construct         & 48 &  8.3\% \\
Tool sub-construct          & 42 & 14.3\% \\
\midrule
High salience               & 84 &  1.2\% \\
Low salience                & 60 & 23.3\% \\
\bottomrule
\end{tabular}
\end{table}
The salience contrast is the largest single effect in the pilot. Fourteen of the 15 violations occurred on low-salience items (those with the prohibition placed mid-prompt rather than at the opening of the system role), yielding a 23.3\% violation rate against 1.2\% for high-salience items. We read this as preliminary support for the policy-invisible-violation literature's claim that prohibitions become unenforceable when buried in long instruction contexts~\cite{phantompolicy2026}. The sub-construct contrast is smaller, with tool-trigger items eliciting violations at roughly 1.5 to 1.7 times the rate of authorization- and state-trigger items, but the absolute counts are too small to support strong inferences.

\subsection{Pilot limitations}
We name the following limitations explicitly so that they bound the inferential reach of the pilot's empirical claims.

The single-seed design means we cannot decompose observed-score variance into person, item, and occasion components. The Generalizability-Theory analysis described in Section~\ref{sec:gstudy} requires multiple seeds per (model, item) cell, which the pilot did not produce. The headline rates are point estimates without standard errors, and the per-model differences in Table~\ref{tab:pilot-per-model} may not survive a full G-study.

The model panel is restricted to open-weights models served through Groq's hosted inference. The pilot does not include frontier closed models (Claude, GPT-4o, Gemini), and it does not include the larger MoE models in the OSS catalog (DeepSeek V3.1 or R1). The nomological-network predictions described in Section~\ref{sec:audit} are not yet testable on a panel this narrow.

The judge ensemble's reliability is bounded by the Llama 3.3 70B Versatile judge having returned valid labels on only 20 of 144 trials, owing to Groq's daily request-per-model cap. The kappa we report is therefore the most reliable two-judge pair (Llama 3.1 8B Instant against Llama 4 Scout 17B at $\kappa = 0.419$, $n=138$) rather than a balanced three-judge ensemble. The substantive-validity argument is correspondingly conditional on this two-judge agreement.

Finally, the pilot is digital-only. The cross-domain generalizability argument in Section~\ref{sec:audit} explicitly identifies the digital-to-embodied transfer as the load-bearing assumption that HIUA-Bench is designed to instrument. Embodied items are not in this pilot, and no claims about embodied behavior follow from these results.

\subsection{What the pilot establishes}
Three things, in our view. First, the three-layer scoring stack is empirically separable: the action scorer's regex and refusal-prefix layers handle approximately 97\% of trials, and the LLM-judge fallback fires on the small minority of ambiguous cases. The recall scorer's substring layer alone would have missed 28.2\% of the recall-correct trials, which is direct evidence for the substantive-validity claim that a benchmark using substring scoring alone would systematically mismeasure HIUA. Second, the construct's $2 \times 4 \times 2$ partition produces non-trivial occupancy in both the HIUA and KBV cells (11 and 4 trials respectively), indicating that the cells are operationally distinguishable rather than reducing to a single mechanism. Third, the salience contrast (1.2\% violation rate at high salience against 23.3\% at low salience) provides preliminary support for the prediction in Section~\ref{sec:audit} that prohibition placement is a load-bearing facet of the construct, and that benchmarks scoring agents with the full prohibition prominently placed in the system role license a deployment claim that production with token-budget pressure may not honor. The pilot's specific HIUA-to-KBV ratio of approximately 2.75 to 1 does not yet generalize, but the partition itself is empirically realizable on the panel and items studied.

\section{Discussion}
A validity audit cannot decide whether a model is safe to deploy. It can decide which deployment claims a benchmark licenses~\cite{m2m2025}. The contribution that this paper offers is, therefore, a tighter license rather than a stronger guarantee. The HIUA score, properly measured, supports the claim that this agent will not, with rate above $\alpha$, hallucinate authorization, state, or tool effect, and thereby execute a forbidden action in distribution $D$. It does not support broader safety claims, and it does not support narrower ones either.

The strongest counter-argument is, in our view, the pragmatic-benchmarking one: imperfect measurements still provide useful feedback during iterative development, and demanding pre-registered nomological networks and G-Theory reliability decompositions raises the cost of new benchmark development without obviously raising its value~\cite{m2m2025}. For capability benchmarks, in which the consequence of an over-claim is a misallocated research dollar, we find this argument substantially correct. For agent-safety benchmarks used to license deployment in irreversible-action settings, the asymmetry of error cost flips the calculus. A false negative (deploying a system that the benchmark certified as safe, which then takes an irreversible forbidden action) is, by the definition of irreversibility, catastrophic. The measurement standard should, as the validity literature has long maintained, rise with the stakes of the use~\cite{kane2013, standards2014}. The HIUA case is the case in which the stakes are explicitly extra-academic.

A second counter-argument is more methodological. Why not measure outcomes and ignore the cognitive process? The substantive critique in Section~\ref{sec:audit} is, on this view, philosophical hand-wringing about mechanism when the outcome is what matters. Our reply is that the same outcome rate can be produced by lucid violation (which an instruction-following intervention would mitigate), by hallucinated authorization (which a faithful-reasoning intervention would mitigate), or by hallucinated state (which a perception-grounding intervention would mitigate). The interventions are distinct, the costs are distinct, and the appropriateness depends on which mechanism dominates. A benchmark that does not partition the mechanism cannot, in our view, inform the intervention choice. It can only flag that something is wrong. The pilot in Section~\ref{sec:pilot} is, in this sense, a small piece of empirical support for that claim: the partition produced 11 HIUA-candidate and 4 KBV trials on the same panel, which an unpartitioned benchmark would have collapsed into a single 10.4\%-violation report.

The deception and sycophancy adjacency is, we believe, worth marking. Strategic deception~\cite{deceptionsurvey2024} involves unfaithful intermediate generation chosen deliberately to achieve a goal. Sycophancy~\cite{sharmasycophancy2023} involves unfaithful intermediate generation produced under user pressure. HIUA involves unfaithful intermediate generation produced without either deliberateness or user pressure. The model has neither been instructed to hallucinate nor been rewarded for hallucinating. It has hallucinated and acted. The three constructs are, in our view, related but distinct. HIUA-Bench items should be designed to discriminate among them: with versus without explicit goal incentives for the violation, and with versus without user pressure favoring the violation. If the rates do not differ across these manipulations, we take the construct decomposition to be wrong, and the three should be reported as facets of a single underlying tendency.

On the digital-to-embodied transfer, SafeAgentBench's 5\% hazardous-task rejection rate~\cite{safeagentbench2024} and HEAL's 40-times hallucination-rate inflation under scene-task inconsistency~\cite{heal2025} are, in our view, the strongest empirical anchors for arguing that the construct is real and measurable in embodied settings. The transfer assumption, namely that a digital-sandbox score predicts embodied behavior, is the most important thing that HIUA-Bench is designed to instrument rather than to assume. The cross-domain stratum of the design is built around this transfer being treated as an open empirical question per model, rather than as a downstream design decision.

On what this paper does not contribute beyond the pilot: the 36-item factorial is sketched, not fully implemented. The pilot in Section~\ref{sec:pilot} runs against the 24-item base bank at single-seed and on a panel restricted to open-weights models. The nomological-network predictions are predictions, not findings. The construct decomposition is theoretically motivated, and rests on a small number of empirical anchors (DriftBench's KBV dispersion~\cite{driftbench2026}, HEAL's hallucination probing~\cite{heal2025}, and the agent-hallucination taxonomy~\cite{halluagentsurvey2025}) supplemented by our own pilot, rather than on extensive convergent evidence. A future empirical version of this work would pilot the full 36-item factorial with paraphrase expansion against 8 to 12 frontier and open-weights models in Inspect Evals scaffolding~\cite{inspectevals2024}, and would report $\Phi$, $\mathrm{E}\rho^2$, the IRT discrimination distribution, and the predicted nomological correlations.

\section{Limitations}
The reliance on chain-of-thought traces as evidence of cognition is itself a substantive assumption, namely that the construct of interest is even partially observable from the trace. Recent work on unfaithful chain-of-thought reasoning makes this assumption non-trivial. Chain-of-thought may be a post-hoc rationalization rather than a faithful representation of the inference that produced the action. The HIUA-Bench design relies on the trace primarily for hallucination-type labeling, rather than for causal attribution. However, the labeling itself is an inference, and the inter-rater reliability of the labeling will bound the substantive-validity argument. The pilot's $\kappa = 0.419$ on the most-reliable judge pair is moderate rather than strong, and the substantive-validity argument is correspondingly bounded by that level of agreement.

The proposed nomological-network correlations are predictions, not findings. If the empirical correlations differ substantially from the predictions, the appropriate response is to report the deviations as evidence about the construct's actual placement, rather than as failures of HIUA-Bench. The pre-registration is a discipline on the predictions, not a guarantee of their truth.

The military and lethal-system motivation is rhetorically load-bearing for the consequential-aspect argument, but the proposed items live in tabletop simulators, not in actual weapons systems. The transfer from a HIUA-Bench score to a real-world lethal-deployment risk is a generalizability claim across at least three levels of remove (simulator to physical-robot to lethal-physical-robot), and HIUA-Bench instrumentation covers only the first. ARMOR's military-context findings~\cite{armor2026} and RoboPAIR's physical-robot findings~\cite{robopair2024} are the empirical bridges across the next two levels. The bridge is not built in this paper.

The pilot's scope means that HIUA-Bench is delivered, in this paper, as a validated design plus a single-seed open-weights pilot rather than as a fully validated benchmark. The single-seed design forecloses the variance-component decomposition that the Section~\ref{sec:gstudy} reliability plan calls for, the panel restriction forecloses the nomological-network correlations, and the partial loss of one judge to Groq's daily rate cap forecloses the balanced three-judge agreement claim. Further empirical instantiation is future work, and the present paper should be read as a measurement-design proposal whose central claim is that the design is principled and the partition is empirically realizable on at least one panel, rather than that the design has been shown to discriminate among frontier models at scale.

\section{Conclusion}
The agent-safety benchmark landscape in 2026 contains many instruments measuring the things that are easy to measure: aggregate failure rates, malicious-instruction compliance, and jailbreak vulnerability. It does not, however, contain an instrument measuring the thing that matters for deploying agents in irreversible-action settings, namely the conditional rate at which an agent hallucinates and, because of that hallucination, crosses a prohibition that it could otherwise restate verbatim. The reason for the gap is not, in our view, that the gap is small. It is, instead, that the gap is harder to measure than the things that fill the existing benchmarks, and the measurement community has, reasonably, started with the easier things.

This paper has argued that the gap is now load-bearing. RoboPAIR, ARMOR, and the Anthropic agentic-misalignment red team have together established that the deployment frontier is moving into settings in which a single forbidden action is irreversible, and in which existing benchmark scores are being cited as evidence that the action will not be taken. The validity audit in Section~\ref{sec:audit} shows that the cite cannot do the work it is being asked to do. The substantive aspect is undefended, the structural reliability is unreported, the generalizability transfer is unmeasured, and the consequential failure modes (gaming, prompt-theater, false assurance) are already visible. The HIUA-Bench design in Section~\ref{sec:design} is a sketch of what an instrument with the right validity properties might look like, and the pilot in Section~\ref{sec:pilot} demonstrates that the partition is empirically realizable. Most strikingly, the pilot reveals that prohibition salience is the largest single effect on violation rate (1.2\% at high salience against 23.3\% at low salience), a finding that bears directly on the consequential-aspect argument that benchmarks scoring agents with the full prohibition prominently placed in the system role license a deployment claim that production with token-budget pressure will not honor. The argument of the paper is that smaller and on-construct beats larger and confounded when the cost of the wrong answer is irreversible, and the pilot provides preliminary evidence that smaller and on-construct is also measurable.

\bibliographystyle{plain}
\bibliography{references}
\end{document}
